% ============================================================
%  数字通信 · 公式记忆手册 —— 可复用模板（lecture-organizer 模式三）
%  用法：复制本文件到 formula-memory-workspace/<章节>/notes.tex，
%        保留导言区与所有 \newcommand，只替换正文。编译：xelatex notes.tex（两遍）。
%  本文件本身可编译（smoke test），xelatex notes.tex 应得 0 错误 0 缺字形。
% ============================================================
\documentclass[11pt,a4paper]{ctexart}
\usepackage{amsmath,amssymb,amsthm,mathtools}
\usepackage{bm}
\usepackage{geometry}
\geometry{a4paper, margin=1.8cm}
\usepackage{enumitem}
\setlist{nosep,leftmargin=1.5em}
\usepackage[most]{tcolorbox}
\tcbuselibrary{breakable, skins}
\usepackage{tabularx}
\usepackage{booktabs}
\usepackage{array}
\usepackage[colorlinks=true, linkcolor=blue!50!black,
            urlcolor=blue!60!black, bookmarksopen=true]{hyperref}

\setlength{\parskip}{3pt}
\linespread{1.06}

% ---------- 缩写宏（本章若需新符号，在下面追加 \newcommand）----------
\newcommand{\xp}{x_{+}}
\newcommand{\xl}{x_l}
\newcommand{\sgn}{\operatorname{sgn}}
\newcommand{\Real}{\operatorname{Re}}
\newcommand{\ihat}[1]{\widehat{#1}}
\newcommand{\inner}[2]{\langle #1,#2\rangle}

% ---------- 答案单元格强调（tabularx 安全：颜色，禁用 \bm/\boldsymbol）----------
\newcommand{\af}[1]{\textcolor{blue!50!black}{\ensuremath{#1}}}
\newcommand{\at}[1]{\textcolor{blue!50!black}{\textbf{#1}}}

% ---------- 标签徽章（公式字段后挂）----------
\newcommand{\bibei}{\colorbox{red!14}{\textcolor{red!70!black}{\,\scriptsize\bfseries\sffamily 必背\,}}}
\newcommand{\ketui}{\colorbox{blue!14}{\textcolor{blue!55!black}{\,\scriptsize\bfseries\sffamily 可推\,}}}
\newcommand{\gainian}{\colorbox{gray!18}{\textcolor{black!55}{\,\scriptsize\bfseries\sffamily 概念\,}}}

% ---------- 文字段标签 ----------
\newcommand{\fld}[2]{\par\noindent\textbf{#1}\quad #2\par\smallskip}

% ---------- 章节大条 ----------
\newcommand{\secbar}[1]{\par\addvspace{12pt}\noindent
  \textcolor{blue!50!black}{\rule[-2pt]{5pt}{16pt}}\hspace{6pt}%
  {\Large\bfseries\textcolor{blue!35!black}{#1}}\par
  \addvspace{2pt}\noindent\textcolor{blue!20}{\rule{\linewidth}{0.6pt}}\par\addvspace{4pt}}

% ---------- 公式族卡片 ----------
\newtcolorbox{familybox}[1]{
  enhanced, breakable, sharp corners,
  colback=blue!1, colframe=blue!45!black, boxrule=0.5pt, leftrule=3pt,
  title={\bfseries\sffamily #1}, fonttitle=\bfseries\sffamily,
  colbacktitle=blue!55!black, coltitle=white,
  left=8pt, right=8pt, top=5pt, bottom=5pt}

\newcommand{\core}[1]{\par\noindent
  {\textcolor{orange!75!black}{\textbf{$\bigstar$ 记忆核心}}}\quad\textit{#1}\par\smallskip}

\begin{document}

% ==================== 标题块（替换为本章） ====================
\begin{center}
{\Huge\bfseries\textcolor{blue!40!black}{数字通信 · 公式记忆手册}}\\[3pt]
{\Large 第 X 章\ \S X.X\ \ <本章标题>}\\[2pt]
{\small\textcolor{gray}{<关键词 1> \,$\cdot$\, <关键词 2> \,$\cdot$\, <关键词 3>}}
\end{center}
\vspace{2pt}
\noindent\textcolor{blue!30}{\rule{\linewidth}{1.2pt}}

% ==================== 一、总览 ====================
\secbar{一、本章公式记忆总览}

<一两句点出本章主线/转换链；列出公式族用普通阿拉伯数字 1 2 3（禁用圆圈数字）；
给三句口诀。>

% ==================== 二、公式族卡片（每族一个 familybox） ====================
\secbar{二、公式族卡片}

\begin{familybox}{公式族 1 · <族名>}
\core{<一句话记忆核心：这条族用什么主公式串起来。>}
\fld{主公式 \bibei}{<端点公式>}          % \bibei 必背 / \ketui 可推 / \gainian 概念
\fld{推导链}{<主公式 $\Rightarrow$ 变形 $\Rightarrow$ 特殊情况>}
\fld{常见结论 \bibei}{<关键结论>}
\fld{易混点}{<本族最易记错处>}
\end{familybox}

% 长公式用 \boxed 居中：
% \par\noindent\textbf{主公式}\ \bibei
% \begin{equation*}\boxed{\;<公式>\;}\end{equation*}

% ==================== 三、必背极简版（红框） ====================
\begin{tcolorbox}[enhanced, breakable, sharp corners,
  colback=red!2, colframe=red!55!black, boxrule=0.6pt,
  title={\bfseries\sffamily 三、必背公式极简版},
  colbacktitle=red!60!black, coltitle=white,
  left=10pt, right=10pt, top=4pt, bottom=4pt]
\renewcommand{\arraystretch}{1.45}
\noindent\begin{tabularx}{\linewidth}{@{}>{\bfseries\color{red!60!black}}c X@{}}
\toprule
\rmfamily\color{black}\textbf{\#} & \textbf{必背公式} \\
\midrule
1 & <必背公式 1> \\
2 & <必背公式 2> \\
\bottomrule
\end{tabularx}
\end{tcolorbox}

% ==================== 四、默写空表（答案列 \dotfill） ====================
\secbar{四、默写空表（遮右栏默写）}

\renewcommand{\arraystretch}{1.7}
\noindent\begin{tabularx}{\linewidth}{@{}c p{7.8cm} X@{}}
\toprule
\textbf{\#} & \textbf{默写项} & \textbf{你的答案} \\
\midrule
1 & <默写项，行内空格用 \underline{\hspace{1cm}}> & \dotfill \\
2 & <默写项> & \dotfill \\
\bottomrule
\end{tabularx}
\renewcommand{\arraystretch}{1.0}

% ==================== 五、默写答案（\clearpage 另起页；答案列 \af/\at） ====================
\clearpage
\secbar{五、默写答案（对答案用）}

\renewcommand{\arraystretch}{1.6}
\noindent\begin{tabularx}{\linewidth}{@{}c p{7.8cm} X@{}}
\toprule
\textbf{\#} & \textbf{默写项} & \textbf{答案} \\
\midrule
1 & <同空表默写项> & \af{X^{*}(f)} 或 \at{<文字答案>} \\
2 & <同空表默写项> & \af{X^{*}(f)} \\
\bottomrule
\end{tabularx}
\renewcommand{\arraystretch}{1.0}

% ==================== 六、易混公式对照（三列） ====================
\secbar{六、易混公式对照}

\renewcommand{\arraystretch}{1.5}
\noindent\begin{tabularx}{\linewidth}{@{}>{\raggedright\arraybackslash}p{4.0cm} >{\raggedright\arraybackslash}X >{\raggedright\arraybackslash}X@{}}
\toprule
\textbf{易混点} & \textbf{正确} & \textbf{常见错误} \\
\midrule
<易混点> & <正确> & <常见错误> \\
\bottomrule
\end{tabularx}
\renewcommand{\arraystretch}{1.0}

% ==================== 七、本章公式总账（灰框速记） ====================
\begin{tcolorbox}[enhanced, breakable, sharp corners,
  colback=gray!4, colframe=gray!55, boxrule=0.5pt,
  title={\bfseries\sffamily 七、本章公式总账（速记）},
  colbacktitle=gray!35!black, coltitle=white,
  left=10pt, right=10pt, top=4pt, bottom=4pt]
\textbf{章节}：<章名>\hfill\textbf{进度}：已完成

\smallskip
\textbf{已整理公式族}：1 <族1>\ 2 <族2>\ ...

\smallskip
\textbf{必背}：见第三节。\quad \textbf{可推}：<...>。

\smallskip
\textbf{易混}：<...>。

\smallskip
\textbf{与其他章节关联}：<点明是哪族延伸，不重复推导>。
\end{tcolorbox}

\vspace{6pt}
\begin{center}\small\textcolor{gray}{数字通信公式记忆手册 \,$\cdot$\, <第 X 章> \,$\cdot$\, 供打印默写}\end{center}

\end{document}
